\section{Glossary}
\label{sec:glossary}

\begin{description}[style=nextline, leftmargin=0pt]
  \item[Backend]
    Implementation of \symbolref{IBackend} for a graphics API (or mock).
  \item[SparseCellGrid]
    Unbounded signed-coordinate simulation storage using non-background
    16$\times$16 chunks and temporary 18$\times$18 halos.
  \item[CanvasView]
    Bounded camera-visible presentation with CPU palette targets, display fade,
    one RGB texture, and one cell-aligned world-space quad.
  \item[Canvas]
    Retained dense compatibility presentation and legacy-import fixture; not a
    second production runtime path.
  \item[Command / token]
    One entry in the render command queue describing a GPU operation.
  \item[CRTP]
    Curiously Recurring Template Pattern; used by \symbolref{Drawable} for
    optional immediate \texttt{DrawImpl}.
  \item[Dirty rect]
    Inclusive AABB of cells/texels that changed; used for sparse visual work
    and subimage uploads.
  \item[Drawable]
    \symbolref{DrawList} entry that emits tokens via \texttt{AppendCommands}
    (preferred) or immediate \texttt{Draw} (fallback).
  \item[DrawList]
    Non-owning ordered list of drawables by layer, with the camera and window,
    for one frame; named \texttt{Scene} until D-E31. It is the
    render-extraction list, not the persistent node hierarchy, a program
    scene, or an immediate GL executor.
  \item[Enrollment]
    Creating GPU resources and receiving handles outside the tight frame loop.
  \item[Fetch set]
    The virtual paths one scene load needs (\texttt{collectSceneFetches}),
    fetched into the guest asset cache before the scene instantiates and held
    there until released, so \texttt{IAssetSource::read} stays synchronous
    (D-E24).
  \item[.ilsc format 2]
    The one UTF-8 JSON scene format every Illumo program reads and writes,
    owned by \texttt{Illumo::Content} (\texttt{IlscCodec},
    \texttt{SceneDocument}). Strict core objects, verbatim namespaced data,
    canonical output; format 1 is refused (D-E19).
  \item[Illumo]
    Engine/host object: owns the services and runs a frame in phases
    (\texttt{beginUpdate}, \texttt{endUpdate}, \texttt{beginRender},
    \texttt{endRender}); it owns no program.
  \item[IllumoContext]
    Frozen non-owning public service bag supplied to the program, its scenes
    and the debug overlay when they start (D-E5/D-E6, D-E31).
  \item[illumo.json]
    The one manifest at the root of every package: id, version, title, kind
    (\texttt{app}, \texttt{content}, or \texttt{mod}), targets, dependencies,
    overlays, and an \texttt{app} or \texttt{mod} section. It replaced
    \texttt{app.json} (D-E22).
  \item[Module]
    A WebAssembly binary (\texttt{.wasm}) that a package's program or a
    compute worker runs. The engine's former \texttt{IModule} registration of
    product screens was removed by D-E31; screens are now program scenes.
  \item[Mount]
    One entry of the host's virtual file tree: a mount point
    (\texttt{/engine}, \texttt{/app}, \texttt{/packages/<id>}, or
    \texttt{/project}) and the layers that supply it. Only \texttt{/project}
    is writable (D-E21).
  \item[Overlay]
    A package layer merged onto \texttt{/app} when the package targets the
    launched application. Layers stack by priority, then dependency order,
    then id; directories union and the topmost file wins. There are no
    whiteouts.
  \item[Package (.ilpk)]
    A unit of content with \texttt{illumo.json} at its root: a loose directory
    or an \texttt{.ilpk} archive, a bounded ZIP subset with stored and deflate
    entries (D-E20). Both forms resolve identically.
  \item[Package root]
    The mount a scene's package-relative references resolve against:
    \texttt{/packages/<id>} under \texttt{/packages}, otherwise the first path
    component (\texttt{/app}, \texttt{/project}, \texttt{/engine},
    \texttt{/local}); see \texttt{scenePackageRoot}.
  \item[Palette]
    CPU table mapping all 256 possible cell-state bytes to RGB target colors.
  \item[PBO]
    Pixel buffer object; async unpack path for texture uploads.
  \item[Program]
    The one WASM program a package runs (D-E31). On the host it is a
    \symbolref{WasmProgram} that \symbolref{RuntimeShell} starts, updates,
    dispatches and asks before the window closes; in the guest a
    \symbolref{GuestProgram} owns what every scene shares and a
    \symbolref{SceneDirector}.
  \item[RuleSet]
    Cellular automaton transition rules (\texttt{nextState} + \texttt{evalCell}).
  \item[Scene (program scene)]
    A \symbolref{ProgramScene}: one screen of a program with its own logic, UI
    and console commands, backed by a \symbolref{SceneInstance} (its content)
    and, with \texttt{HostRender}, its own host render world (D-R30). It
    starts once, enters and leaves any number of times, and stops once; a
    scene left is kept, frozen, until the program releases it (D-E31). The
    per-frame drawable list is \symbolref{DrawList}.
  \item[SceneDirector]
    Owns a program's scenes by name and switches between them at frame
    boundaries (a cut, or a cover once the program's transition cover is
    drawn). Only the active scene updates, receives input and draws; each
    scene keeps its own camera.
  \item[SceneGraph]
    Persistent scene-owned hierarchy addressed by graph-local generational
    handles. It caches world transforms and emits visible, enabled borrowed
    render attachments as one \symbolref{DrawList} drawable.
  \item[SceneInstance]
    \texttt{Illumo::Content}'s live scene: it owns one document state, a
    SceneGraph, the attachments that draw each node and their asset
    references. Per-node edits change only what they touch; stable ids are
    graph node names.
  \item[VFS]
    Virtual file system: the host's \texttt{VirtualFileSystem}, an immutable
    mount table of absolute, case-sensitive \texttt{/}-separated paths over
    directory and archive backends. Guests list, stat and read it through
    file protocol v2; nothing it returns is a host path.
\end{description}
